\section{Research question and precise contribution}

The computational objective is a complete unknot recognizer with a
quasipolynomial worst-case bound in the crossing number. The maintained
ProveIt implementation combines exact filters, group-theoretic methods,
normal-surface certificates, and a complete fallback. Many local operations
are now substantially faster than their first implementations. The central
research issue is whether the number and encoded size of the states visited
can be bounded without expanding exponentially large surfaces or returning
to an exhaustive search over all normal-coordinate sectors.

This continuation works on two interfaces where the existing research
packages leave concrete opportunities. In the normal-coordinate search,
syntactically allowed quadrilateral types need not be geometrically
possible. A certificate of the entire active support makes it possible to
remove these spurious choices in one operation and define the remaining
ambiguity intrinsically. In geometric descent, many currently legal
Pachner reductions have disjoint tetrahedral supports. They can be verified
and applied together, avoiding repeated emission and validation of almost
identical triangulations. Vertex-link peeling then requires careful
collective scoring: disjoint tetrahedral supports can still share a global
vertex, so their peeled scores need not add.

\subsection{What is established here}

The mathematical results fall into five groups.
\begin{enumerate}
\item A primal--dual description of the maximal active support of a
nonnegative homogeneous matching cone. A single aggregate dual witness
certifies every universally zero coordinate. Combining this with a
positive objective gives a complete support certificate, independently
checkable without rerunning the search procedure. The underlying
polyhedral principle is classical strict complementarity; the contribution
is its concrete certificate use at this repository interface.
\item An ambiguity parameter for restricted quadrilateral search, together
with a finite search-tree bound, and a negative result: on an unrestricted
one-vertex anchored root that has a positive-Euler witness, the ambiguity
rank is at least $t-1$. More generally it is at least $t-v(T)$,
and at least $16\max\{1,n\}-2$ on the maintained pulling source. The obstruction is explicit and comes from local
directions invisible to all face-arc equations.
\item A simultaneous relative-boundary replacement theorem for disjoint
$3$--$2$ regions. The conflict graph of distinct legal degree-three central
edges has maximum degree at most nine. A greedy maximal batch therefore
contains at least a tenth of the current candidates. The cocycle transport
and unpeeled cell changes commute and add.
\item A deterministic data structure that maintains vertex-link minima
under repeated local replacements. Cached thirteen-record prefixes give
constant-operation hypothetical singleton scoring, while balanced-tree
updates replenish these prefixes after commits. Collective scoring accounts
for shared vertices. A realized solid-torus example proves that singleton
peeled gains cannot simply be added.
\item An exact maximum-weight-matching reduction for the best peeled-Euler
subbatch of any fixed disjoint family. A minimum-omission optimum omits
at most one move per relevant interior vertex. On the maintained pulling
source, which has exactly two interior poles, a linear arithmetic-time
selector omits at most two moves. Thus the packed optimum retains at least
$\max\{0,\lceil M/10\rceil-2\}$ current candidates.
\end{enumerate}

All statements are scoped to their actual inputs. A positive-Euler relaxed
matching vector may violate the quadrilateral constraints. A legal Pachner
move preserves the ambient manifold, but coherent surfaces transported
through it can change topology. A positive total Euler characteristic is
therefore not an unknot certificate. The supplied code retains the
repository's independently verified component-disc query as the relevant
terminal geometric test.

\subsection{Position relative to the latest incoming work}

The source revision inspected is \basecommit. Its \texttt{fast/} and
\texttt{synthesis/} files used here were verified against their Git blob
identities. The three relevant archives in \texttt{docs/incoming/} were
also verified byte-for-byte against the repository's recorded blobs:
the dual-certificates report, the support-sensitive-certificates report,
and the geometric-transport report, all dated 9 October 2026
\cite{priorDual,priorSupport,priorTransport}. They are prior work, not new
deliverables of this continuation.

The dual report introduced independently replayable LP objective
certificates and forced-type propagation. The support report sharpened
component-incidence bounds and primitive ray certification. The transport
report established the exact five-height Pachner score, coherent cocycle
transport, and the thirteen-record lemma for a \emph{fixed} scoring pass.
It explicitly left that last lemma unimplemented and warned that a short
prefix must be replenished after repeated deletions. The present dynamic
construction supplies that missing maintenance mechanism and connects it
to simultaneous geometric replacements.

The package contains a pinned supporting snapshot, a small overlay for
the new modules, an integration patch, exact proof fixtures, retained
measurements, and reproduction scripts. It is prepared for review and
integration; no remote repository modification is part of this task.

\subsection{The literature does not remove the global gap}

Lackenby's technical announcement describes a hierarchy strategy with a
quasipolynomial running-time target \cite{LackenbyTalk}. The particular
Oxford handout gives $2^{O((\log n)^3)}$, which is
$n^{O((\log n)^2)}$. This should not be silently replaced by the sharper
$n^{O(\log n)}$ expression sometimes quoted in public descriptions.
Both are quasipolynomial, but the exponent matters in a rigorous cost
ledger.

The July 2026 preprint gives a substantial hierarchy-based recognition
algorithm and new certificate results \cite{Lackenby2026}. Its discussion
of running time explicitly leaves algorithmic precision and speed-up
for further work. Proposition~10.2 bounds the length of a partial
hierarchy by a linear function of quadrilateral handle complexity; it
does not establish the missing polylogarithmic effective-depth statement
for the code developed here.

The normal-coordinate part also has important established precedents.
Burton--Ozlen's tree traversal uses LP feasibility and greedy support
minimization to obtain a connected positive-Euler surface
\cite{BurtonOzlen}. Burton's admissible-face theory identifies compatibility
with the geometry of nonnegative normal cones \cite{BurtonFaces}. The new
support certificate must be understood against those results. We neither
claim that LP support reduction is new nor infer a polynomial recognition
algorithm from the polynomial solvability of a single LP.
